\section{Prolongation monodromique et calcul à profondeur prescrite}
\label{sec:generacion-monodromica-certificada-rev7}
\index{monodromie nonadique}\index{génération structurelle}

La construction part de l'état enrichi produit par
\(\mathrm{APP}\to\mathrm{TRIT}\to\mathrm{TPK}\) : catalogue actif de
\(468\) émissions, projection ternaire, action \(D_3\), calibre
d'orientation, registre de transitions de rang plein, caractères de
clôture, de propagation et d'auto-échelle, cylindres ternaires, phase nonadique et
ombre de Paley--Witt. Le registre reconstruit de manière unique les
relèvements \(L_0,L_1\) avant d'ouvrir la carte archimédienne. Aucun
chiffre de \(\pi\), de \(e\) ou de \(\varphi\) n'est reçu en entrée du générateur.

\subsection{Première prolongation et frontière coinductive}

Le recensement exhaustif sélectionne les mots
\[
 w_6^\pi=010211,\qquad
 w_6^e=201101,\qquad
 w_6^\varphi=121200.
\]
Avec des vecteurs lignes sur \(\mathbb F_3^6\), les deux élévations reconstruites
sont
\[
L_0=
\begin{pmatrix}
2&2&2&1&2&1\\
2&1&2&2&1&1\\
1&0&0&1&0&2\\
0&0&1&1&1&0\\
2&2&2&0&2&1\\
2&1&2&1&0&0
\end{pmatrix},
\qquad
L_1=
\begin{pmatrix}
2&1&2&0&2&2\\
1&2&1&1&2&0\\
1&2&1&1&1&2\\
0&1&0&0&2&1\\
2&0&2&1&0&1\\
0&2&2&2&1&1
\end{pmatrix}.
\]
Soient \(B^{\mathsf C}\), \(B^{\mathsf P}\) et \(B^{\mathsf A}\) les
biographies structurelles de clôture, de propagation et d'auto-échelle qui sont
publiées ci-dessous, décomposées en blocs
\(b_0^\chi|\cdots|b_4^\chi\). Formons
\[
 \begin{aligned}
 X_0&=(b_0^{\mathsf C},b_1^{\mathsf C},
       b_0^{\mathsf P},b_1^{\mathsf P},
       b_0^{\mathsf A},b_1^{\mathsf A})^T,\\
 Y_0&=(b_1^{\mathsf C},b_2^{\mathsf C},
       b_1^{\mathsf P},b_2^{\mathsf P},
       b_1^{\mathsf A},b_2^{\mathsf A})^T,
 \end{aligned}
\]
et, de manière analogue, \(X_1,Y_1\) avec les transitions
\(b_2\to b_3\) et \(b_3\to b_4\). Chaque bloc s'obtient par la projection
ternaire de l'état produit après la mise à jour complète
\[
 x_{j+1}^{\chi}
 =\mathcal U_{9j+9}\circ\cdots\circ\mathcal U_{9j+1}(x_j^{\chi}),
 \qquad
 b_j^{\chi}=\Pi_6(x_j^{\chi}),
 \qquad
 \mathcal U_t=\mathsf{Upd}_t\circ\mathsf{Tra}_t\circ\mathsf{Sel}_t.
\]
Ainsi, le registre qui précède les relèvements est publié
intégralement par les quatre matrices
\begingroup\scriptsize
\begin{equation}
\begin{aligned}
X_0&=\begin{pmatrix}
0&1&0&2&1&1\\0&1&2&2&2&2\\2&0&1&1&0&1\\
1&2&1&2&2&1\\1&2&1&2&0&0\\1&1&2&2&0&2
\end{pmatrix},&
Y_0&=\begin{pmatrix}
0&1&2&2&2&2\\0&1&0&2&1&1\\1&2&1&2&2&1\\
1&0&2&0&1&1\\1&1&2&2&0&2\\1&2&1&0&2&0
\end{pmatrix},\\[1ex]
X_1&=\begin{pmatrix}
0&1&0&2&1&1\\0&0&2&1&1&1\\1&0&2&0&1&1\\
0&1&2&2&2&2\\1&2&1&0&2&0\\0&1&0&2&1&0
\end{pmatrix},&
Y_1&=\begin{pmatrix}
0&0&2&1&1&1\\1&1&0&2&2&1\\0&1&2&2&2&2\\
1&0&2&0&1&1\\0&1&0&2&1&0\\0&1&0&2&0&0
\end{pmatrix}.
\end{aligned}
\label{eq:registro-transiciones-tpk-explicito-rev10}
\end{equation}
\endgroup
Leurs lignes sont des sorties successives de l'endomorphisme TPK sur les trois germes
sélectionnés par les invariants de fibre ; ce ne sont ni des coefficients imposés à
\(L_0,L_1\) ni des données tirées des développements archimédiens.

\begin{theorem}[Reconstruction unique des relèvements et du calendrier]
\label{thm:reconstruccion-lifts-calendario-rev10}
Dans \(\mathbb F_3\), on a
\[
 \det X_0=1,
 \qquad \det X_1=-1,
 \qquad \operatorname{rank}(I_6\otimes X_i)=36.
\]
Par conséquent, les systèmes
\[
 X_0L_0=Y_0,
 \qquad X_1L_1=Y_1
\]
possèdent les solutions uniques \(L_i=X_i^{-1}Y_i\), qui sont exactement les
matrices imprimées ci-dessus. Parmi les \(16\) calendriers binaires de longueur
\(4\), \(0011\) est le seul qui conserve simultanément les trois
biographies. Les \(144\) substitutions unitaires possibles dans les deux matrices
détruisent cette triple compatibilité.
\end{theorem}

\begin{proof}
Les déterminants et les rangs sont calculés par élimination exacte modulo \(3\).
L'inversibilité de \(X_i\) donne l'existence et l'unicité de \(L_i\) ; les
\(16\) calendriers et les \(144=2\cdot6\cdot6\cdot2\) perturbations forment
deux recensements finis exhaustifs. Toutes leurs entrées sont des mots, des transitions
et des résidus de l'état enrichi antérieur à l'évaluation réelle. Les noms
conventionnels des constantes n'interviennent dans aucun de ces calculs.
\end{proof}

La suite
\[
u_1=u_0L_0,\qquad u_2=u_1L_0,\qquad
u_3=u_2L_1,\qquad u_4=u_3L_1
\]
produit
\[
\begin{aligned}
w_{30}^{\pi}&=010211|012222|010211|002111|110221,\\
w_{30}^{e}&=201101|121221|102011|012222|102011,\\
w_{30}^{\varphi}&=121200|112202|121020|010210|010200.
\end{aligned}
\]
La matrice
\[
R_{36}=
\begin{pmatrix}
222220\\021222\\102011
\end{pmatrix}
\]
est la première frontière coinductive. Les matrices \(L_0,L_1\) sont des résultats
de la reconstruction interne de
\cref{thm:reconstruccion-lifts-calendario-rev10} et s'appliquent avant d'ouvrir
les témoins décimaux. Le domaine du théorème est l'unique chaîne génératrice
\(\mathrm{APP}\to\mathrm{TRIT}\to\mathrm{TPK}\to\) état enrichi.
APP apporte ses deux feuillets arithmétiques ; TRIT oriente le régime ; TPK construit
et conserve les champs de route, de retenue, de frontière et de mémoire sur lesquels agissent
les relèvements. Demander que la projection visible d'APP contienne à elle
seule l'état terminal confond deux domaines consécutifs de cette même
chaîne et ne transforme pas les relèvements en données introduites à la main.

\subsection{Sélection finie et construction des caractères}

Soient \(\mathcal C\) le multiensemble de germes, \(\mathcal U\) le catalogue
d'émissions et \(W=\Psi(\mathcal U)\subset\mathbb F_3^6\), avec
\[
 \mathcal C\xrightarrow{T_{54}^{\rm seed}}\mathcal U
 \xrightarrow{\ \Psi\ }W\lhook\joinrel\longrightarrow\mathbb F_3^6,
 \qquad
 \Psi(U_6)=((-u_j)\bmod3)_{j=1}^6.
\]
L'énumération exacte donne
\[
 \sum_{U\in\mathcal U}\operatorname{mult}(U)=104\,976,\qquad
 |\mathcal U|=468,\qquad |W|=243,\qquad|\mathbb F_3^6|=729.
\]
Ces quatre nombres comptent des objets typés distincts.

L'action \(D_3\) sur \(W\) produit \(43\) orbites. L'orbite de clôture
est la seule dont les six orientations possèdent cinq élévations, une masse \(1008\)
et des multiplicités \(432+4\cdot144\). Parmi les orbites radicales unitaires,
le recensement orienté de la clôture détermine de manière unique la propagation et le
recensement \((2,2,2)\) détermine de manière unique l'auto-échelle. Le calibre
\(\operatorname{diag}_c=6\), avec orientation ES, fixe un mot dans chaque
orbite. Par conséquent, le triplet
\[
 (w_6^\pi,w_6^e,w_6^\varphi)
 =(010211,201101,121200)
\]
s'obtient avant de construire tout intervalle archimédien.

Les trois caractères s'obtiennent également à partir de la structure sélectionnée. Pour
la clôture, la matrice d'incidence
\[
 A_W=\begin{pmatrix}
 0&1&1&1&1&1\\
 1&0&1&2&2&1\\
 1&1&0&1&2&2\\
 1&2&1&0&1&2\\
 1&2&2&1&0&1\\
 1&1&2&2&1&0
 \end{pmatrix}
 \quad\text{satisfait}\quad A_W^2=2I_6=-I_6
\]
et fixe l'inversion orientée. Sur la pentafibre, le symétriseur qui
oublie l'étiquette de l'élévation est
\[
 P_5=\frac15\mathbf1\mathbf1^T.
\]
La conservation \(\mathbf1^T\lambda=1\) et l'invariance
\(P_5\lambda=\lambda\) imposent \(\lambda=(1/5)\mathbf1\), donc
\(t_1=1/5\). Les quatre
compositions par
\[
 t\star s=\frac{t+s}{1-ts}
\]
donnent \(t_4=120/119\). L'équation de retour à la diagonale,
\[
 \frac{120/119-1/q}{1+(120/119)/q}=1,
\]
possède l'unique solution entière positive \(q=239\). L'identité gaussienne
\begin{equation}
 \label{eq:generacion-monodromica-identidad-gaussiana-rev7}
 (5+i)^{16}(239-i)^4=-681386607803576157184
\end{equation}
reconnaît ensuite le demi-tour. Ainsi, l'expression
\[
 16\arctan(1/5)-4\arctan(1/239)
\]
est l'évaluation exacte du caractère de clôture, non son sélecteur.

Pour la propagation, la composition des chemins et l'unité conservée donnent
\[
 P_{s+t}=P_sP_t,\qquad P_0=1,\qquad [t]P_t=1.
\]
Si \(P_t=\sum_{n\geq0}a_nt^n\), la comparaison du coefficient de \(s^nt\)
impose
\[
 (n+1)a_{n+1}=a_n,\qquad a_0=1,
\]
et, par récurrence, \(a_n=1/n!\). L'égalité \(P'_t=P_t\) et la valeur
\(P_1=e\) sont des reconnaissances analytiques ultérieures de cette loi formelle.

Pour l'auto-échelle, le calendrier \((+,-,0)\) induit
\[
 F_{\rm av}=\begin{pmatrix}0&1\\1&1\end{pmatrix}.
\]
Son unique rayon positif invariant \(v=(1,x)^T\) satisfait
\[
 F_{\rm av}v=xv,\qquad x^2-x-1=0,\qquad x>0.
\]
Les quotients consécutifs de Fibonacci fournissent des encadrements rationnels
emboîtés pour \(x\) ; l'équation \(x=1+1/x\) est la coordonnée scalaire de
cette équation propre, non une condition imposée de l'extérieur.

\begin{proposition}[Séparation entre sélection, évaluation et reconnaissance]
\label{prop:generacion-monodromica-separacion-causal-rev7}
Pour chaque canal \(\chi\), la composition causale est
\[
 \widehat X
 \xrightarrow{\operatorname{Sel}_{\rm TPK}}
 X_\chi
 \xrightarrow{\operatorname{ev}_{\mathbb R}}
 \mathbb R
 \xrightarrow{\operatorname{Rec}}
 \{\pi,\varphi,e,\alpha\}.
\]
\(\operatorname{Sel}_{\rm TPK}\) utilise le recensement, l'orbite, le calibre,
les relèvements, la retenue et le caractère interne. L'évaluation contracte
les cylindres compatibles. La reconnaissance attribue le nom conventionnel
à la valeur déjà construite. Aucun préfixe cible n'intervient dans les deux premières
flèches.
\end{proposition}

\begin{proof}
L'unicité des trois orbites se décide au moyen de prédicats entiers du
catalogue. Les paramètres \(5,4,239\), la récurrence
\((n+1)a_{n+1}=a_n\) et la matrice \(F_{\rm av}\) se déduisent de ces orbites et
des lois d'unité, de composition et d'incidence. La coordonnée de
\(\alpha\) provient ensuite du triplet interne au moyen du registre signé
\(U_{12}\), du sceau \(K\) et de la retenue dodécaphasique. Les bornes rationnelles sont
calculées uniquement pour évaluer et certifier la section déjà fermée. La
comparaison décimale s'effectue ensuite ; elle ne peut donc pas agir comme
entrée causale.
\end{proof}

\subsection{Extension des survivants et certification ultérieure}

Soit \(\mathcal H_m^\chi\) l'ensemble des histoires enrichies admissibles
du canal \(\chi\) à la profondeur \(m\), et soit
\(p_{n,m}:\mathcal H_n^\chi\to\mathcal H_m^\chi\) la troncature. On définit
\[
 \operatorname{Surv}_{m}^{(N),\chi}
 =p_{N,m}(\mathcal H_N^\chi),\qquad
 \operatorname{Surv}_{m}^{\chi}
 =\bigcap_{N\ge m}\operatorname{Surv}_{m}^{(N),\chi},
\]
et la section numérique par
\[
 x_\chi\in X_\chi:=
 \varprojlim_m\operatorname{Surv}_m^\chi.
\]
Le lecteur et l'état ne sont pas identifiés : une présentation
\((x_\chi,L_\chi)\) publie la valeur, tandis que l'état conserve aussi
l'orientation, la retenue, la frontière, le feuillet, l'incidence exceptionnelle et la généalogie.
À la frontière R36, \(\chi_{\rm HMT}\) et le résidu interne
\(\varepsilon_K^\chi\) sont déjà des coordonnées de l'état. La division APP et le
sélecteur de phase produisent
\[
 d_K^\chi=\operatorname{Dig}_{729}(\varepsilon_K^\chi),\qquad
 \varepsilon_{K+1}^\chi=729\varepsilon_K^\chi-d_K^\chi,qquad
 N_{K+1}^{\operatorname{Ext}}=729N_K+d_K^\chi,
\]
et la sélection, le transport et la mise à jour
\[
 \operatorname{Sel}_K^\chi(x)
 :=\operatorname{Sel}\!\left(
 x,M_{\rm ph}(\rho_9(K+1))\right),
 \qquad
 \operatorname{Tra}_K^\chi(y,d):=\operatorname{Lift}(T_d)(y),
 \qquad
 \operatorname{Upd}_K^\chi:=\operatorname{Rec}_K^\chi\circ
 \operatorname{Mem}_K^\chi,
\]
\[
 \mathcal U_K^\chi(x)
 :=\operatorname{Upd}_K^\chi\!\left(
  \operatorname{Tra}_K^\chi\!\left(
   \operatorname{Sel}_K^\chi(x),d_K^\chi
  \right)\right)
\]
dépend uniquement de l'état précédent. La relation complète
\(\operatorname{Ext}\) peut être multivaluée ; la sous-relation
\(\operatorname{Ext}^{\chi,\rightarrow}\), ancrée dans l'état R36 qui
porte déjà le caractère, est le graphe de \(\mathcal U_K^\chi\). Elle fixe donc
une section globale compatible sans affirmer l'unitarité de chaque fibre locale de
\(\operatorname{Ext}\). Ce n'est pas une intersection rétrospective définie par
\(P_{\rm ar}\), et elle ne lit ni chiffres futurs ni encadrements.

Si \(N_K\) est le préfixe ternaire et \(L_K=6K\), ses \(729\) sous-cylindres sont
\[
 I_{K,u}=
 \left[
 \frac{729N_K+u}{3^{L_K+6}},
 \frac{729N_K+u+1}{3^{L_K+6}}
 \right),
 \qquad 0\le u<729.
\]
Pour l'encadrement rationnel exact \([\ell_{\chi,K},h_{\chi,K}]\) construit
par la récurrence interne du caractère \(\chi\),
\[
\begin{aligned}
j_{\min}
 &=\max\!\left(729N_K,
   \left\lfloor\ell_{\chi,K}3^{L_K+6}\right\rfloor\right),\\
j_{\max}
 &=\min\!\left(729N_K+728,
   \left\lceil h_{\chi,K}3^{L_K+6}\right\rceil-1\right).
\end{aligned}
\]
La vérification archimédienne ultérieure est
\[
 j_{\min}=j_{\max}=N_{K+1}^{\operatorname{Ext}}.
\]
Elle certifie que la partition est exhaustive et localise le survivant déjà
produit par \(\operatorname{Ext}^{\chi,\rightarrow}\) ; elle ne le sélectionne pas.
L'encadrement n'est pas lu dans un développement décimal : il est calculé
à partir de la pentafibre et du retour unitaire, de la récurrence de propagation
ou de l'incidence \(F_{\rm av}\), selon le canal. Par conséquent,
l'intersection est un certificat de publication, non une règle de descendance.

\begin{theorem}[Génération coinductive corrélée à profondeur arbitraire]
\label{thm:generacion-monodromica-conjunta-rev7}
Dans la structure enrichie TPK construite dans les théorèmes précédents, les
trois caractères dérivés et la fermeture dodécaphasique construisent à profondeur
arbitraire la famille corrélée de sections
\(x_\pi,x_\varphi,x_e,x_\alpha\) qui satisfait :
\begin{enumerate}
 \item el barrido
 \[
 104\,976\longrightarrow468\longrightarrow243
 \]
 sépare les germes, les émissions et les mots visibles ;
 \item los levantamientos
 \[
 w_6\xrightarrow{L_0}w_{12}\xrightarrow{L_0}w_{18}
 \xrightarrow{L_1}w_{24}\xrightarrow{L_1}w_{30}
 \]
 précèdent la première frontière coinductive \(R_{36}\) ;
 \item à chaque niveau,
 \(\operatorname{Ext}^{\chi,\rightarrow}\) produit d'abord le fils
 compatible et sa troncature coïncide avec l'état du niveau précédent ; les
 \(729\) sous-cylindres numériques forment une partition ordonnée et l'encadrement
 rationnel localise ensuite le fils déjà engendré ;
 \item la phase satisfait \(g(K+9)=g(K)\), tandis que le bloc, l'indice de
 cylindre, le préfixe et l'ombre de Witt changent ;
 \item les lecteurs publient, respectivement, \(\pi\), \(\varphi\), \(e\)
 et \(\alpha\), et la règle d'extension n'a pas de profondeur terminale.
\end{enumerate}
\end{theorem}

La portée de l'énoncé est coinductive et de profondeur arbitraire. Pour les
quatre canaux, l'unicité de la section forward et sa compatibilité sous
troncature proviennent de la mise à jour de l'état. On n'affirme pas que toute
fibre locale de \(\operatorname{Ext}\) soit unitaire. Les exécutions à toute
profondeur finie ne sont que des témoins de régression et de publication.

\begin{proof}
La sélection finie de
\cref{prop:generacion-monodromica-separacion-causal-rev7} conserve les cinq
régions de l'orbite de clôture avant d'ouvrir la carte réelle. La pente
primitive \(1/5\), la composition des quatre branches accompagnatrices et la
condition de retour unitaire imposent, dans cet ordre, \(120/119\) et \(q=239\).
L'identité
\eqref{eq:generacion-monodromica-identidad-gaussiana-rev7} reconnaît alors
le premier demi-tour positif, de sorte que l'évaluateur ultérieur est
\[
 \pi_{\rm HMT}
 =16\arctan(1/5)-4\arctan(1/239).
\]
Les arctangentes sont évaluées par des séries alternées avec des bornes rationnelles.

En propagation, la loi \(P_{s+t}=P_sP_t\), l'unité \(P_0=1\) et la
normalisation linéaire imposent d'abord \((n+1)a_{n+1}=a_n\). Ce n'est qu'ensuite que l'on
reconnaît \(P'_t=P_t\) et
\[
 e_{\rm HMT}=\sum_{n\ge0}\frac1{n!}.
\]
En auto-échelle, le calendrier produit d'abord
\[
 F_{\rm av}=\begin{pmatrix}0&1\\1&1\end{pmatrix}
\]
et son rayon de Perron. L'équation \(x=1+1/x\) est sa coordonnée scalaire
ultérieure.

La mise à jour qui ne lit que l'état précédent produit dans chaque canal le
fils compatible à partir du caractère et du résidu internes présents. Par
récurrence, deux sections forward ayant le même état R36 coïncident à toute
profondeur. Dans les trois canaux à lecteur réel développés ici,
l'arithmétique rationnelle produit ensuite des intervalles de largeur arbitrairement
petite : la partition exacte des \(729\) sous-cylindres vérifie que le fils
déjà engendré se trouve dans la cellule décimale publiée. La fermeture signée
\(U_{12}\to K\) prolonge la quatrième coordonnée et les carrés de troncature
de ses deux voies internes commutent, de sorte que
\(\alpha_A=\alpha_B=\alpha_{\rm HMT}\). Les exécutions à mille et dix mille
chiffres sont des projections finies, non le critère du théorème.
\end{proof}

L'extension à une section de la limite inverse est régie par
\cref{thm:generacion-supervivencia-pi-e-phi-rev6}. Le théorème précédent
établit, pour \(\pi,\varphi,e,\alpha\), la non-vacuité, la compatibilité et
l'unicité globale de chaque section forward exigées par ce principe
abstrait, sans postuler l'unitarité porte par porte de
\(\operatorname{Ext}\). Cette formulation distingue la règle coinductive de
profondeur arbitraire de ses évaluations finies à mille et dix mille chiffres,
sans modifier le retour de phase ni le transport de mémoire.

\begin{corollary}[Profondeur arbitraire et objet décimal infini]
\label{cor:generacion-nonadica-profundidad-arbitraria-rev10}
Pour tout \(N\in\mathbb N_{>0}\), la même construction publie un unique
préfixe décimal canonique \(d_N^\chi\) de longueur \(N\), avec
\(d_{N+1}^\chi|_N=d_N^\chi\), pour
\(\chi\in\{\pi,\varphi,e,\alpha\}\). La famille compatible
\((d_N^\chi)_{N\ge1}\) détermine le développement infini dans la limite
inverse. Mille, dix mille ou un million de chiffres sont des exécutions finies du
même algorithme paramétré ; aucune d'entre elles ne limite la portée du
théorème.
\end{corollary}

\begin{proof}
La mise à jour \(\operatorname{Upd}_\chi\) produit d'abord le cylindre
compatible. L'encadrement rationnel du caractère a une largeur arbitrairement
petite et certifie ensuite que ce cylindre rencontre un unique membre de
la partition publiée des \(729\) sous-cylindres. Étant donné \(N\), il existe \(m\) avec
\(3^{-6m}<10^{-N}\) ; le cylindre déjà produit est alors contenu dans une
cellule décimale canonique de longueur \(N\). La compatibilité de troncature
identifie ces cellules comme les coordonnées finies d'une section forward.
\end{proof}

\begin{theorem}[Génération coinductive corrélée et double projection]
\label{thm:estado-coinductivo-doble-proyeccion-rev10}
Soit \(\mathscr S_\infty\) le produit fibré des quatre sections
forward corrélées de profondeur arbitraire engendré par
\[
 \mathrm{APP}\longrightarrow\mathrm{TRIT}\longrightarrow\mathrm{TPK}
 \longrightarrow R_{36}\longrightarrow G_9.
\]
De \(\mathscr S_\infty\) partent deux publications typées :
\[
 \mathscr S_\infty\xrightarrow{\mathcal P_{\rm Arq}}\mathbb R^4,
 \qquad
 \mathscr S_\infty\xrightarrow{\mathcal P_{\rm Exc}}
 W_{12}\longrightarrow G_{24}\longrightarrow\Lambda_{24}
 \longrightarrow V^\natural
 \longrightarrow\operatorname{Aut}(V^\natural)=\mathbb M.
\]
La première contracte les cylindres compatibles et publie de manière corrélée
\((\pi_{\rm HMT},\varphi_{\rm HMT},e_{\rm HMT},\alpha_{\rm HMT})\) ; la
seconde conserve l'incidence de frontière et publie la chaîne exceptionnelle.
Par conséquent, \(\mathbb R\) est l'ombre archimédienne de la première
projection et non une entrée qui sélectionne la généalogie. Machin,
Perron--Fibonacci, le semi-groupe exponentiel et les dénominations
conventionnelles agissent après la sortie HMT.
\end{theorem}

\begin{proof}
\cref{thm:generacion-monodromica-conjunta-rev7} et
\cref{cor:generacion-nonadica-profundidad-arbitraria-rev10} construisent la
famille compatible avant l'évaluation. La fermeture dodécaphasique publie
\(\alpha_{\rm HMT}\) à partir de son triplet R36 corrélé et du registre
\(U_{12}\to K\) ; la voie
\(E_{21/22}\) conserve son domaine analytique propre. La projection
archimédienne oublie l'incidence en contractant les cylindres, tandis que la
projection exceptionnelle la transporte par les opérateurs d'entrelacement Witt--Golay--
Leech jusqu'à \(V^\natural\). Toutes deux proviennent du même antécédent et aucune
ne peut agir rétroactivement comme générateur de celui-ci.
\end{proof}

\subsection{Régions de clôture}

Les cinq régions ne sont pas cinq valeurs de \(\pi\), mais cinq généalogies
ayant une projection visible commune. Leur multiplicité est
\[
 432+4\cdot144=1008,
\]
et leurs rôles sont
\[
 3_{++}+1_{--}+1_{\perp}.
\]
L'identité régionale se vérifie par le quadruplet
\[
 (U_6,E_{\rm sig},\text{multiplicité},\text{rôle}),
\]
non seulement par le recensement \(3/1/1\). En particulier,
\[
\begin{array}{c|c|c|c}
U_6&E_{\rm sig}&\text{multiplicité}&\text{rôle}\\ \hline
501|614|498|169|272|272&555555&432&\perp\\
810|923|870|169|272|272&339555&144&++\\
810|983|810|169|272|272&393555&144&++\\
870|923|810|169|272|272&933555&144&++\\
870|923|810|418|674|521&933717&144&--
\end{array}
\]
La table fixe conjointement le mot décimal enrichi, la signature d'émission,
la multiplicité et le rôle. En particulier, l'ancrage de multiplicité (432) est
transversal et la signature (933717) est le retour muté ; le recensement (3/1/1)
ne détermine pas à lui seul ces noms.

\subsection{Retour de phase et changement de mémoire}

À la neuvième phase,
\[
 B_9=(100100\mid020112\mid010122)
\]
admet trois orientations qui conservent l'axe \(\varphi\) et l'agrégat
\(\pi+e\) :
\[
\begin{aligned}
 &(101222\mid112221\mid112002),\\
 &(102221\mid111222\mid112002),\\
 &(111221\mid102222\mid112002).
\end{aligned}
\]
Les trois survivent à une frontière et conservent
\[
 u^{\varphi}=112002,
 \qquad u^\pi+u^e=210110.
\]
La frontière suivante sélectionne
\[
 B_{10}=(101222\mid112221\mid112002).
\]
Par conséquent,
\[
 \#\operatorname{Surv}_1(B_9)=3,
 \qquad
 \#\operatorname{Surv}_2(B_9)=1.
\]
L'indice dix occupe de nouveau la phase un, mais non l'état précédent. C'est
la monodromie : retour de phase avec transformation de mémoire.

En termes de l'état complet
\[
 \widehat x_K=(B_K,C_K,p_K,c_K,\Sigma_K,W_K,g(K)),
\]
le tour est l'holonomie unique
\[
 \begin{gathered}
 \Gamma_9=\operatorname{Hol}_{C_9}(\nabla_{\rm TPK}),
 \qquad g(K+9)=g(K),\\
 q_{\rm mem}\Gamma_9=q_{\rm mem}+1,
 \qquad \Gamma_9\widehat x_K\ne\widehat x_K.
 \end{gathered}
\]
Dans une réalisation isométrique, la même continuité satisfait
\[
 U_{\Gamma_9}J=JT+\eta,\qquad
 I-T^*T=\eta^*\eta.
\]
La retenue se compose par
\[
 \operatorname{Carr}(\gamma_2\gamma_1)
 =\operatorname{Carr}(\gamma_2)H(\gamma_1)
 +Y(\gamma_2)\operatorname{Carr}(\gamma_1),
\]
et la télescopie conserve le terminal fini :
\[
 S_0=\sum_{r=0}^{L-1}M_9^{*r}D_9^*D_9M_9^r
     +M_9^{*L}S_0M_9^L.
\]
Les fenêtres \(27\), \(54\) et \(108\) sont les concaténations
\[
 \Gamma_{27}=\Gamma_9^3,\qquad
 \Gamma_{54}=\Gamma_9^6,\qquad
 \Gamma_{108}=\Gamma_9^{12};
\]
transportent une mémoire accumulée de \(3\), \(6\) et \(12\) tours,
respectivement. Elles ne constituent pas des réinitialisations de l'algorithme.

\subsection{Recensement fini des transitions et génération décimale}

L'exécution à profondeur mille produit, par canal,
\[
396\ \text{transitions},\quad
2376\ \text{trits},\quad
43\ \text{tours complets},\quad
387\ \text{paires }(K,K+9),
\]
et \(1000\) chiffres décimaux. À chaque transition, la mise à jour
\(\operatorname{Ext}^{\chi,\rightarrow}\) produit le survivant et le
balayage exhaustif ultérieur des \(729\) sous-cylindres certifie sa
publication. Les témoins externes sont consultés après la génération et
utilisés uniquement pour la reconnaissance.

L'indépendance de la cible s'obtient parce que l'application génératrice reçoit
uniquement le catalogue fini des \(468\) émissions, les matrices
\(L_0,L_1,A_W,F_{\rm av}\) et l'arithmétique rationnelle. Aucune de ses fonctions
de transition ne contient de préfixes de \(\pi\), de \(e\), de \(\varphi\) ni une
opération qui évalue ces constantes. Les encadrements se déduisent des
caractères construits et n'agissent que dans l'évaluation ultérieure.

Après les 396 étapes de prolongation et de certification, le cylindre final
\[
 \left[\frac{N}{3^{2376}},\frac{N+1}{3^{2376}}\right)
\]
est contenu dans une unique cellule décimale de profondeur \(10^{-1000}\).
Les mille chiffres sont donc un témoin fini de la loi paramétrable, non
sa profondeur maximale.

L'argument uniforme construit un cylindre pour toute profondeur
finie prescrite. Le recensement des \(396\) étapes ne vérifie que l'instance à
mille chiffres : il maintient séparées la génération causale et la comparaison décimale
ultérieure, et ne remplace pas le quantificateur de
\cref{thm:generacion-monodromica-conjunta-rev7}.
